\section{Discussion}
In this work, we analyzed how attention weights contribute to fairness and accuracy of a predictive model. To do so, we proposed an attribution method that leverages the  attention mechanism and showed the effectiveness of this approach on both tabular and text data. Using this interpretable attribution framework we then introduced a post-processing bias mitigation strategy based on attention weight manipulation. We validated the proposed framework by conducting experiments with different baselines, as well as fairness metrics, and different data modalities.

Although our work can have a positive impact in allowing to reason about fairness and accuracy of models and  reduce their bias, it can also have negative societal consequences if used unethically. For instance, it has been previously shown that interpretability frameworks can be used as a means for fairwashing which is when malicious users generate fake explanations for their unfair decisions to justify them \cite{pmlr-v119-anders20a}. In addition, previously it has been shown that interpratability frameworks are vulnerable against adversarial attacks \cite{slack2020fooling}. We acknowledge that our framework may also be targeted by malicious users for malicious intent that can manipulate attention weights to either generate fake explanations or unfair outcomes. An important future direction can be to analyze and improve robustness of our framework along with others. 

%Another possible direction for future work is to design a method that can identify the right threshold value as a flag that a feature is problematic, or being more cautious about fairness-accuracy trade-off and keeping features useful for accuracy and even up-weight them even if they contribute to unfairness. 
%In this work, we considered all the features whose removal would result in a lower bias than the original model as problematic features that should be removed without giving much attention to accuracy. However, there is room for improvement to design methods that can learn the optimum threshold value for removal of these features based on the desired accuracy vs fairness trade-off. This can be considered as a limitation of the current work that we plan to address in the future. 

